% 本文件由 LaTeX 排版 Agent 根据有效 translation.json 自动生成。
% author=Councils; work_id=3804; title=《4世纪的刚格拉主教会议》

\chapter{《4世纪的刚格拉主教会议》}

% structure: p0001 source=h1 target=omitted reason=page-title-already-rendered-by-parent

% structure: p0002 source=h2 target=section reason=relative-heading-rank; base=h2

\section{背景}

关于刚格拉主教会议，除了从其本身的会议通告中所了解的以外，我们所知甚少。对此自然会提出三个重大问题。\par

1. 其日期为何？\par

2. 其所谴责的欧斯塔修是何人？\par

3. 谁是其主席？\par

我将简要地向读者提供关于这些问题的要点。\par

1. 关于日期，毫无疑问是在尼西亚会议之后、君士坦丁堡第一届大公会议之前，即在325年至381年之间。索克拉特似乎将其定在约365年；但索佐门则提前了约二十年。另一方面，雷米·塞利耶与其其他陈述不一致，似乎根据圣巴西略的书信论证真正的日期晚于376年。还有另一种由巴莱里尼兄弟提出的理论，基于主持的优西比乌是凯撒勒雅的优西比乌这一假设，因此他们将日期定在362年至370年之间。与福克斯先生同意此说，并与帕奇一起定在358年，并大胆地补充说：\Quote{这无疑是该会议的年分。}但在几乎所有的古代典集汇编中，刚格拉典集均先于安提约基雅典集，布隆德尔和蒂勒蒙支持了这一观点，或许我可以称其为传统日期。\par

2. 似乎没有什么合理理由怀疑被谴责者，即名叫欧斯塔修的人，正是著名的塞巴斯特主教。这一点可以从索佐门和索克拉特两人那里得知，并由圣巴西略的一封信函间接证实。此外，欧斯塔修的塞巴斯特教区在亚美尼亚，而会议正将其信函致给亚美尼亚的主教们。考虑到这一切，赫费勒主教的话似乎并不过分，他写道：\Quote{在这种情况下，巴罗尼奥、杜宾等人（没有古代任何单一证言支持）的主张，即这里指的是另一个欧斯塔修，或可能是隐修士欧塔克图斯，不值得认真考虑，尽管蒂勒蒙并未表示反对。}\par

至于欧斯塔修在刚格拉主教会议谴责后放弃其奇特服饰和其他怪癖的说法，索佐门仅作为传闻提及。\par

3. 关于谁担任主席，似乎可以相当肯定其名为优西比乌——如果索佐门确实是指\Quote{君士坦丁堡的优西比乌}，那是一个错误，但他名字是对的。在会议通告的开头，优西比乌名列第一，由于刚格拉和亚美尼亚属于凯撒勒雅管辖范围，自然可以假设所提及的优西比乌是该省的总主教，但必须记住，卡帕多细雅的优西比乌直到362年才被任命为主教，比福克斯先生认为他主持刚格拉会议的时间晚了四年。希腊文本中列出了十三位主教的名字。\par

拉丁文译本添加了其他名字，例如科尔多瓦的奥西乌，一些拉丁文作者断言他作为教宗的\Emph{\OriginalTerm{特使}{à latere}}主持会议，例如巴罗尼奥和比尼乌斯。赫费勒否认这一点，并说：\Quote{在刚格拉主教会议时，奥西乌无疑已经去世。}但学术界并非一贯如此认为，卡夫认为奥西乌的主教任期长达七十年，结束于361年，并且（基于同一观点）帕奇认为奥西乌可能在358年返回西班牙途中参加了该会议，正如我说过的，福克斯先生同意这一观点。此外，到六世纪初，刚格拉主教会议显然在罗马被视为在教宗权威下举行；教宗西玛克在504年罗马主教会议上明确如此声明。（见第七、第八条典集注释）\par

剩下的只需补充说明，\WorkTitle{会议文集}提到一位名叫迪乌斯的人担任会议主席。巴莱里尼兄弟建议应读作\OriginalTerm{（希腊文）}{\GreekText{Βίος}}，即优西比乌的缩写。福克斯先生建议迪乌斯是\Quote{可能是优西比乌的前任迪亚尼乌斯}。莱特福特将巴勒斯坦的优西比乌的主教任期定在313年至337年之间；并指出卡帕多细亚凯撒勒雅的优西比乌的主教任期直到362年才开始，因此读者将看到巨大的年代学困难。\par

由于所有提出的新日期都或多或少涉及矛盾，我将这些典集放在它们通常的位置，即新凯撒勒雅会议与安提约基雅会议之间，并暂不确定日期。\par

% structure: p0015 source=h2 target=section reason=relative-heading-rank; base=h2

\section{会议通告}

优西比乌、埃利安、欧根尼、奥林匹乌、比提尼库、额我略、斐勒图、帕普斯、欧拉里、希帕提乌、普罗埃西乌、巴西尔和巴苏斯，在刚格拉圣主教会议聚集，致我们最尊敬的主内同工、亚美尼亚的弟兄们，愿在主内安康。\par

\Emph{既然}最神圣的主教会议为处理教会若干必要事务而聚集于\Place{冈格拉}，在审查有关\Person{欧斯大特}的事项时，发现这些追随\Person{欧斯大特}的人做了许多不法之事，故被迫作出决议，并迅速公诸于众，以纠正他所行的一切错谬。因为他们极端憎恶婚姻，并主张凡处于婚姻状态者对天主毫无希望，导致许多已婚妇女离弃丈夫，丈夫离弃妻子，之后因不能自制而陷于奸淫，从而因这种原则而蒙羞。此外，他们还煽动人们离开天主的殿宇和教会，藐视教会及其成员，私设聚会和集会，散布与教会及其礼仪相悖的教义和其他事物；穿着奇装异服，破坏普遍服装习惯；将教会初熟之物（自始即归于教会者）分给自己及其追随者，视之为圣人；奴隶离弃主人，并因自己的奇装异服而对主人无礼；妇女不顾体统，穿男装而非女装，以为可因此成义；许多人假借虔诚之名剪除妇女天然的发辫；他们在主日禁食，藐视这自由日的神圣，却在教会规定的斋日轻蔑地进食；其中有些人厌恶食肉；也不容忍在已婚者家中祈祷，反而藐视这些祈祷，并常常拒绝领受在已婚者家中举行的献礼；他们轻视已婚的司铎，拒绝接触其职务；他们谴责纪念殉道者的礼仪以及在其中聚会或服务的人，也谴责那些没有放弃所有财富的富人，认为他们对天主毫无希望；还有许多无法尽述的事。因为他们每个人都背离了教会的教规，各自制定了近乎隔离的法则；他们之间也没有共同的意见，而是各人随己意设想，提出主张，致使教会受辱，自身遭毁灭。\par

因此，聚集于\Place{冈格拉}的神圣主教会议被迫据此谴责他们，并宣布将他们逐出教会；但若他们悔改，并诅咒这些谬误教义中的每一条，则应可恢复教籍。所以，神圣主教会议特别列出了他们在被接纳之前必须诅咒的一切。若有人不服从上述法令，应作为异端者处以绝罚，逐出教会；主教们亦应对所有与他们同流合污者遵循同样规则。\par

% structure: p0019 source=h2 target=section reason=relative-heading-rank; base=h2

\section{教规}

% structure: p0020 source=h3 target=subsection reason=relative-heading-rank; base=h2

\subsection{第一条教规}

若有人谴责婚姻，或厌恶、谴责一位虔诚的信女与自己的丈夫同寝，认为她不能进入（天国），则处以绝罚。\par

% structure: p0022 source=h3 target=subsection reason=relative-heading-rank; base=h2

\subsection{第二条教规}

若有人谴责那食肉（即未放血、未祭偶像、未被勒死的肉）且信仰虔诚的人，认为此人因食肉而毫无（救恩的）希望，则处以绝罚。\par

% structure: p0024 source=h3 target=subsection reason=relative-heading-rank; base=h2

\subsection{第三条教规}

若有人假借虔诚之名，教导奴隶藐视主人、逃避服务，而不以善意和尊敬服事自己的主人，则处以绝罚。\par

\BibleRef{弟前 6:1}\newline{}这些经文同样被\Person{巴尔萨蒙}和\Person{佐纳拉斯}引用。\par

% structure: p0027 source=h3 target=subsection reason=relative-heading-rank; base=h2

\subsection{第四条教规}

若有人声称，对于已婚的司铎，不可领受他所奉献的祭品，则处以绝罚。\par

% structure: p0029 source=h3 target=subsection reason=relative-heading-rank; base=h2

\subsection{第五条教规}

若有人教导藐视天主的殿及其中的集会，则处以绝罚。\par

% structure: p0031 source=h3 target=subsection reason=relative-heading-rank; base=h2

\subsection{第六条教规}

若有人在教会之外私自集会，并藐视教规，擅自举行教会职务，而司铎未经主教同意拒绝许可，则处以绝罚。\par

\EditorialNote{这两条教规（第五和第六条）禁止秘密聚会及其仪式。从冈格拉主教会议信函的第二条已可看出，欧斯大特派出于灵性骄傲，自视为纯洁神圣，与其余会众分离，避开公共崇拜，举行私密仪式。从信函的第九、十、十一条可知，欧斯大特派尤其避免在已婚圣职人员主礼时参加公共礼仪。我们或许可以从第六条教规的措辞\Quote{没有司铎在场，或未经主教同意}推断，没有任何司铎参与他们的私人仪式；但更可能的是，欧斯大特派并不拒绝司铎职本身，而是厌恶已婚圣职人员，他们有自己的未婚圣职人员，并由这些人在其分离仪式中主礼。上述教规的措辞与此假设并不矛盾，因为\Quote{\OriginalTerm{未经主教的同意}{\GreekText{κατὰ} \GreekText{γνώμην} \GreekText{τοῦ} \GreekText{ἐπισκόπου}}}这一附加语表明，主持欧斯大特派仪式的派别司铎并未获得当地主教的许可。这就是希腊注释家\Person{巴尔萨蒙}等以及\Person{范埃斯彭}对这条教规的解释。}\par

% structure: p0034 source=h3 target=subsection reason=relative-heading-rank; base=h2

\subsection{第七条教规}

若有人擅自取用献给教会的果实，或在教会之外施予这些果实，而未获主教或负责此类事务之人的同意，且拒绝服从其决定，则处以绝罚。\par

% structure: p0036 source=h3 target=subsection reason=relative-heading-rank; base=h2

\subsection{第八条教规}

若有人（除主教或负责管理善款之人外）给予或收受收入，则给予者与收受者均处以绝罚。\par

% structure: p0038 source=h3 target=subsection reason=relative-heading-rank; base=h2

\subsection{第九条教规}

若有人因厌恶婚姻而守童贞或节欲，而非出于童贞本身的美善与神圣，则处以绝罚。\par

\Person{佐纳拉斯}写道：童贞最是美好，节欲亦然，但仅当我们因其本身及其所带来的圣化而遵循它们时才是美好的。但若有人因厌恶婚姻为不洁而守童贞，保持独身，并因认为与女性交往和婚姻本身为邪恶而禁绝之，则此教规对其处以绝罚。\par

这一条规收录于《教会法大全》（\WorkTitle{Corpus Juris Canonici}）中格拉提安（Gratian）的《教令集》（\WorkTitle{Decretum}）第一部分，第三十题第五章，亦见于第三十一题第九章。\par

% structure: p0042 source=h3 target=subsection reason=relative-heading-rank; base=h2

\subsection{第十条规}

若有人因主的缘故而守童贞生活，却傲慢地对待已婚者，此人应受绝罚。\par

% structure: p0044 source=h3 target=subsection reason=relative-heading-rank; base=h2

\subsection{第十一条规}

若有人藐视那些因信德举办爱筵、为荣耀主而邀请弟兄们的人，并且因轻视所行之事而不愿接受这些邀请，此人应受绝罚。\par

% structure: p0046 source=h3 target=subsection reason=relative-heading-rank; base=h2

\subsection{第十二条规}

若有人假借苦修之名，穿戴\OriginalTerm{披风}{peribolæum}，并自以为因此获得义德，从而藐视那些虔敬地穿着\OriginalTerm{短披肩}{berus}及使用其他普通惯常服饰的人，此人应受绝罚。\par

\EditorialNote{\OriginalTerm{短披肩}{\GreekText{βήροι}}（即 \OriginalTerm{短披肩}{lacernæ}）是男子穿在束腰外衣外面的普通外衣；而\OriginalTerm{粗布披风}{\GreekText{περιβόλαια}}是哲学家所穿的粗布披风，以示他们蔑视一切奢侈。苏格拉底（\WorkTitle{教会史}，ii. 43）和冈格拉主教会议信函第三条提到，塞巴斯特的欧斯塔修斯（Eustathius of Sebaste）穿着哲学家的披风。但本条规绝非绝对禁止隐修士穿着特殊服装，因为会议在此谴责的并非独特的服装，而是对其价值骄傲而迷信的过高估计。}\par

% structure: p0049 source=h3 target=subsection reason=relative-heading-rank; base=h2

\subsection{第十三条规}

若任何妇女假借苦修之名，更换服装，不穿女子惯常的衣服，而穿上男子的衣服，她应受绝罚。\par

\BibleRef{申 22:5} 因此，整个古代教会最严格地禁止此事。此类换装先前主要用于戏剧目的，或出于娇柔、放荡、助长不贞等。欧斯塔修斯派出于完全相反和极端苦修的理由，曾建议妇女穿上男性的、可能即隐修士的服装，以表明她们作为圣者已不再有性别之分；但教会出于苦修的理由，也禁止这种换装，尤其是当它与迷信和清教徒式的骄傲相结合时。\par

% structure: p0052 source=h3 target=subsection reason=relative-heading-rank; base=h2

\subsection{第十四条规}

若任何妇女离弃她的丈夫，并因厌恶婚姻而决意离开他，她应受绝罚。\par

\EditorialNote{范埃斯彭（Van Espen）注：欧斯塔修斯派似乎主要不赞同婚姻的运用，并以保持节制的借口，诱导已婚妇女将其视为非法而予以避免，并离开她们的丈夫，至少在夫妻关系上分居。希腊注释者如此理解本条规；因为从未有人指控欧斯塔修斯派劝导他人解除婚姻的约束。}\par

% structure: p0055 source=h3 target=subsection reason=relative-heading-rank; base=h2

\subsection{第十五条规}

若有人离弃自己的子女，不养育他们，也不尽其所能以相称的虔敬教养他们，反而假借苦修之名忽视他们，此人应受绝罚。\par

\EditorialNote{本次会议的教父们在此教导：父母有责任和本分照顾子女的肉身，并尽其所能塑造他们实践虔敬。父母对子女的这种照顾应优先于任何私人的宗教修习。在这方面应参阅圣方济各·沙雷氏（St. Francis de Sales）的书信（第四卷第三十二封）。}\par

% structure: p0058 source=h3 target=subsection reason=relative-heading-rank; base=h2

\subsection{第十六条规}

若任何子女假借虔敬之名离弃其父母——尤其是当父母为信徒时——并以尊重虔敬胜过父母为由，不向父母表示应有的尊敬，他们应受绝罚。\par

\EditorialNote{佐纳拉斯（Zonaras）指出，\Quote{尤其是}一词表明此项义务是普世性的。所有注释者均在此引用玛窦福音第十五章，主在那里谈到犹太人假借虔敬的托词欺骗父母，使天主的法律失效。}\par

% structure: p0061 source=h3 target=subsection reason=relative-heading-rank; base=h2

\subsection{第十七条规}

若任何妇女假借苦修之名剪去头发——这头发是天主赐予她作为顺服的提醒，从而如同废除了顺服的规范——她应受绝罚。\par

\BibleRef{格前 11:10}\par

% structure: p0064 source=h3 target=subsection reason=relative-heading-rank; base=h2

\subsection{第十八条规}

若有人假借苦修之名在星期日守斋，他应受绝罚。\par

\EditorialNote{欧斯塔修斯规定主日为斋日，然而由于基督在这一天从坟墓中复活，并将人性从罪恶中解救出来，我们应以喜悦的感恩奉献给天主度过此日。但守斋带有悲伤和哀痛的含义。因此，凡在星期日守斋者均受绝罚的惩罚。}\par

% structure: p0067 source=h3 target=subsection reason=relative-heading-rank; base=h2

\subsection{第十九条规}

若任何苦修者，没有身体上的必要，却傲慢无礼，并因自以为对此事有完全的理解而忽视教会普遍规定并遵守的斋戒，他应受绝罚。\par

% structure: p0069 source=h3 target=subsection reason=relative-heading-rank; base=h2

\subsection{第二十条规}

若有人出于傲慢的心态，谴责并憎恶〔纪念〕殉道者的集会、或在那里举行的礼仪、以及他们的纪念活动，他应受绝罚。\par

% structure: p0071 source=h3 target=subsection reason=relative-heading-rank; base=h2

\subsection{结语}

我们写下这些，并非要断绝那些愿意依照圣经在天主教会内度苦修生活的人，而是要断绝那些假借苦修之名而行傲慢之事的人：他们既高抬自己，超过那些生活简朴的人，又引入违反圣经及教会规条的新奇事物。我们确实钦佩与谦卑相伴的童贞，尊重与虔敬和庄重相伴的节制，赞扬以谦卑之心离弃世俗事务；同时，我们尊重婚姻的神圣共融，不藐视以正直和行善所享有的财富；我们赞赏衣着上的朴素和节俭——这只出于对身体并非过分讲究的注意——而摒弃放纵和娇柔过度的服饰；我们尊敬天主的殿宇，并将其中举行的集会视为神圣和有益的，不将宗教局限于房屋内，而是尊敬每一个奉天主之名建造的地方；我们赞同在教会内聚集以谋求共同的益处；并祝福弟兄们按照教会传统对穷人所行的极大慈善；简而言之，我们希望在教会内遵守一切由圣经和宗徒传统所传授的。\par

\EditorialNote{这份结语在古代摘要中缺失；虽然它在狄奥尼修斯·埃克西古斯（Dionysius Exiguus）和伊西多尔·梅卡托尔（Isidore Mercator）的译本中出现在第二十条规之后，但并未被编为一条规。此外，在安提约基雅的若望（John of Antioch）的汇编和佛提乌（Photius）的《教规汇编》（Nomocanon）中，规条数目据说为二十条。只有希腊注释者将其编为第二十一条规，但其真实性毫无争议。}\par

\SourceRecord{Translated by Henry Percival.；Nicene and Post-Nicene Fathers, Second Series；Vol. 14.；Edited by Philip Schaff and Henry Wace.；Buffalo, NY: Christian Literature Publishing Co.,；1900.；<http://www.newadvent.org/fathers/3804.htm>.}
